\documentclass[10.5pt,a4paper]{article}
\usepackage[margin=2.1cm,top=2.2cm,bottom=2.2cm]{geometry}
\usepackage{fontspec}
\usepackage{amsmath,amssymb,mathtools}
\setmainfont{Latin Modern Roman}
\newfontfamily\headfont{Inter}[UprightFont=* SemiBold,BoldFont=* Bold]
\newfontfamily\labelfont{Inter}[UprightFont=* Medium]
\usepackage{xcolor}
\usepackage{booktabs,tabularx,array,colortbl}
\usepackage[most]{tcolorbox}
\usepackage{enumitem}
\usepackage{titlesec}
\usepackage{fancyhdr}
\usepackage{ragged2e}
\usepackage{needspace}
\usepackage[hidelinks]{hyperref}

\definecolor{ink}{HTML}{1F2328}
\definecolor{muted}{HTML}{5C6470}
\definecolor{felt}{HTML}{1E6B4A}
\definecolor{feltsoft}{HTML}{EAF4EE}
\definecolor{line}{HTML}{CBD5CF}
\definecolor{warnsoft}{HTML}{FBF3E4}
\definecolor{warn}{HTML}{9A6A12}
\color{ink}

\raggedbottom
\setlength{\parindent}{0pt}
\setlength{\parskip}{4pt}
\RaggedRight
\hyphenpenalty=10000
\sloppy

\titleformat{\section}{\headfont\Large\color{felt}}{\thesection}{0.6em}{}[\vspace{-2pt}{\color{line}\rule{\linewidth}{0.6pt}}]
\titleformat{\subsection}{\headfont\normalsize}{\thesubsection}{0.5em}{}
\titlespacing*{\section}{0pt}{14pt}{6pt}
\titlespacing*{\subsection}{0pt}{10pt}{3pt}

\pagestyle{fancy}
\fancyhf{}
\renewcommand{\headrulewidth}{0pt}
\fancyfoot[L]{\labelfont\footnotesize\color{muted} Bender · Hoja de fórmulas · Parte 1: modelo probabilístico}
\fancyfoot[R]{\labelfont\footnotesize\color{muted}\thepage}

% Caja de fórmula
\newtcolorbox{formbox}[2][]{enhanced,before upper={\RaggedRight},colback=white,colframe=line,boxrule=0.6pt,arc=3pt,
  left=8pt,right=8pt,top=6pt,bottom=6pt,
  title={\headfont\small #2},coltitle=felt,colbacktitle=feltsoft,
  attach boxed title to top left={yshift=-2.2mm,xshift=6pt},
  boxed title style={boxrule=0pt,arc=2pt,colback=feltsoft},#1}
\newtcolorbox{notebox}{enhanced,breakable,before upper={\RaggedRight},colback=warnsoft,colframe=warnsoft,boxrule=0pt,arc=3pt,left=8pt,right=8pt,top=5pt,bottom=5pt}

% Tabla de parámetros
\newenvironment{params}{\par\smallskip\renewcommand{\arraystretch}{1.25}\small
  \tabularx{\linewidth}{@{}>{\raggedright\arraybackslash}p{2.9cm}>{\raggedright\arraybackslash}X@{}}
  \toprule {\labelfont Parámetro} & {\labelfont Qué es y de dónde sale}\\ \midrule}
  {\bottomrule\endtabularx\par}
\newcommand{\uso}[1]{\par\smallskip{\labelfont\small\color{felt} Cómo usarla.}\ {\small #1}}
\newcommand{\ej}[1]{\par\smallskip{\labelfont\small\color{muted} Ejemplo.}\ {\small #1}}
\newcommand{\code}[1]{{\ttfamily\small #1}}

\let\oldsection\section
\renewcommand{\section}{\needspace{16\baselineskip}\oldsection}
\begin{document}

{\headfont\color{felt}\Huge Bender\par}
\vspace{2pt}
{\headfont\Large Hoja de fórmulas · Parte 1: modelo probabilístico\par}
\vspace{4pt}
{\color{muted}Las fórmulas exactas con las que se implementan la Capa 1 (cerebro sólido), la Capa 2 (lectura de rivales) y la explotación. La Parte 2 (RL: bandido y ajuste de parámetros) va en una hoja aparte.\par}
\vspace{6pt}

\begin{notebox}\small
\textbf{Cómo leer esta hoja.} Cada caja trae la fórmula, una tabla con cada parámetro y una línea de cómo usarla. Los valores marcados como \emph{provisional} son puntos de partida: se miden con los bots del kit y el servidor de pruebas, y se afinan en la Capa 3. Los valores marcados como \emph{calibrado} ya se midieron con los 6 bots de entrenamiento del kit (paso 4); la versión vigente de cada número vive en \code{params.py}. Todas las cantidades de fichas se miden \textbf{desde la decisión actual}: lo que ya se puso en el pozo no cuenta (costo hundido), así que $EV(\text{retirarse}) = 0$ siempre.
\end{notebox}

\section*{0 · Notación común}

\begin{params}
$P$ & (Sola, sin paréntesis.) Pozo en el momento de decidir, \textbf{incluidas} todas las apuestas de esta ronda (también la del rival que hay que igualar). \\
$C$ & Fichas que Bender tiene que agregar para igualar. \\
$B$ & Tamaño de una apuesta cuando nadie apostó todavía en la ronda. \\
$A,\ A'$ & En una apuesta o subida: $A$ = fichas que \textbf{agrega Bender}; $A'$ = fichas que agrega el rival que la paga. \\
$e$ & Equity de Bender: fracción del pozo que se lleva en promedio si la mano llega al showdown. \\
$e'$ & Equity de Bender \textbf{si le pagan} (contra el rango que sigue después de pagar). \\
$c_i$ & Probabilidad de que el rival $i$ \textbf{pague} (o siga) ante una apuesta. \\
$f$ & \textbf{Fold equity}: probabilidad de que \textbf{todos} se retiren. Solo se usa para esto. \\
$x$ & Tamaño de una apuesta de Bender, en fracción del pozo (definición general en §1.8). \\
$z$ & Tamaño de una apuesta \textbf{observada de un rival}. Se calcula igual que $x$; cambia solo quién apuesta (§2.7). \\
$h$ & Un combo concreto de dos cartas del rival, por ejemplo Q$\spadesuit$Q$\heartsuit$. \\
$w(h)$ & Peso del combo $h$ en el rango del rival, entre 0 y 1. \\
$t$ & Tipo de rival: nit, TAG, station, maníaco o aleatorio. \\
$\pi_t$ & Probabilidad (creencia actual) de que el rival sea de tipo $t$. \\
$\mathrm{bb}$ & Ciega grande: 100 fichas (ciegas fijas 50/100 todo el torneo). \\
$P(\cdot)$ & Con paréntesis es una \textbf{probabilidad}, por ejemplo $P(a\mid t)$. No confundir con el pozo $P$. \\
\end{params}

% =====================================================================
\clearpage
\section{Capa 1 · Cerebro sólido}

\subsection{Combos y bloqueadores}
\begin{formbox}{Combos de cada clase de mano}
\[
\#\text{pareja}=\binom{4}{2}=6 \qquad \#\text{suited}=4 \qquad \#\text{offsuit}=4\cdot 3=12 \qquad \text{total}=1326
\]
Con cartas muertas (las de Bender y el board), se cuentan solo los combos \textbf{vivos}:
\[
\#\text{pareja viva}=\binom{n_r}{2} \qquad \#\text{suited viva}=q \qquad \#\text{offsuit viva}=n_{r_1}\,n_{r_2}-q
\]
\begin{params}
$n_r$ & Cuántas cartas del rango $r$ (por ejemplo, ases) siguen vivas: 4 menos las que ve Bender. \\
$q$ & Palos en los que \textbf{las dos} cartas siguen vivas (cada uno da un combo suited). \\
\end{params}
\uso{Antes de cualquier cálculo de rango, se borran los combos que usan una carta muerta.}
\ej{Bender tiene A$\spadesuit$J$\heartsuit$: AA vivos $=\binom{3}{2}=3$. Para AK: $n_A=3$, $n_K=4$, $q=3$ ($\heartsuit,\diamondsuit,\clubsuit$) $\Rightarrow$ AKs $=3$, AKo $=12-3=9$.}
\end{formbox}

\subsection{Equity por Monte Carlo}
\begin{formbox}{Equity y su error}
\[
e \approx \frac{W + \sum_{j\in\text{empates}} 1/k_j}{N}
\qquad\qquad
\sigma_e=\sqrt{\frac{e(1-e)}{N}}
\qquad\qquad
N \ge \frac{e(1-e)}{\varepsilon^2}
\]
\begin{params}
$N$ & Número de simulaciones (repartos al azar del resto del board y de las manos rivales). \\
$W$ & Simulaciones en que Bender gana solo. \\
$k_j$ & En la simulación $j$ con empate, cuántos jugadores empatan (Bender cobra $1/k_j$). Heads-up es $1/2$. \\
$\sigma_e$ & Error típico de la estimación. \\
$\varepsilon$ & Error máximo aceptado. \\
\end{params}
\uso{Para $\varepsilon = 1\,\%$ en el peor caso ($e=0{,}5$) basta $N = 2500$. Cuándo dejar de simular lo decide el reloj adaptativo de abajo.}
\end{formbox}

\begin{formbox}{Cuándo dejar de simular (reloj adaptativo)}
\[
\text{parar si}\quad N\ge N_{\min}\ \text{ y }\ \Big(\sigma_e\le\varepsilon_{\text{obj}}\ \ \text{o}\ \ |e-e_{\text{ref}}|\ge k\,\sigma_e\Big)
\qquad\text{o si}\quad t\ge T_{\max}
\]
\[
T_{\max}=\min\!\Big(T_{\text{tope}},\ \eta\cdot\frac{T_{\text{rest}}}{D_{\text{rest}}}\Big)
\qquad
D_{\text{rest}}=2{,}5\cdot(50-\text{manos jugadas en la mesa})
\]
\begin{params}
$N_{\min}$ & Mínimo de simulaciones antes de poder parar: 2000. \\
$\varepsilon_{\text{obj}}$ & Precisión suficiente: 0{,}005 (medio punto de equity). \\
$e_{\text{ref}}$ & El umbral que importa para esta decisión, el más cercano a $e$ entre: $e_{\min}$ de §1.6 (si hay que igualar), 0{,}25 y 0{,}50 (límites de §1.11). \\
$k$ & Cuántos $\sigma_e$ de distancia al umbral bastan para que más simulaciones no cambien la decisión: 3. \\
$t$ & Tiempo ya usado en esta decisión, revisado cada 64 simulaciones. \\
$T_{\text{tope}}$ & Tope absoluto por decisión: 1{,}5 s (el límite de la arena es 5 s). \\
$T_{\text{rest}}$ & Segundos que le quedan a la mesa: $360$ menos el tiempo desde la primera mano de la mesa. \\
$D_{\text{rest}}$ & Decisiones de Bender que faltan, estimadas con unas 2{,}5 por mano. \\
$\eta$ & Fracción de su parte justa del reloj que Bender se permite usar: 0{,}5. \\
\end{params}
\uso{Cada mesa se corta a los 6 minutos aunque no lleguen a 50 manos. Con unos 7 s por mano para toda la mesa, Bender no puede gastar 2 s fijos por decisión. Medido con el evaluador en C (\code{phevaluator}): mediana de 5 ms y p95 de 17 ms por decisión; el tope casi nunca actúa.}
\end{formbox}

\subsection{Equity contra un rango}
\begin{formbox}{Promedio ponderado por el rango}
\[
e_R=\frac{\sum_h w(h)\,e(h)}{\sum_h w(h)}
\qquad\qquad
P(\text{sortear } h)=\frac{w(h)}{\sum_{h'} w(h')}
\]
\begin{params}
$e(h)$ & Equity de Bender si el rival tuviera exactamente el combo $h$. \\
$w(h)$ & Peso del combo vivo $h$ en el rango (§2.1). \\
\end{params}
\uso{No se calcula combo por combo: en cada simulación de Monte Carlo se \textbf{sortea} el combo del rival con probabilidad proporcional a $w(h)$. Con varios rivales se sortea un combo de cada rango y se descarta la simulación si dos combos comparten carta.}
\ej{A$\spadesuit$J$\heartsuit$ contra \code{QQ+, AK} (27 combos vivos): $e_R=6{,}768/27\approx 25\,\%$, contra $63{,}7\,\%$ frente a una mano al azar.}
\end{formbox}

\subsection{Stack efectivo y modo de juego}
\begin{formbox}{Stack efectivo en ciegas grandes}
\[
S_{\text{ef}}=\frac{\min\!\big(S^{0}_{\text{Bender}},\ \max_i S^{0}_i\big)}{\mathrm{bb}}
\qquad
\text{modo}=\begin{cases}
\text{push/fold} & S_{\text{ef}}\le 15\\
\text{intermedio} & 15< S_{\text{ef}}\le 25\\
\text{profundo} & S_{\text{ef}}>25
\end{cases}
\]
\begin{params}
$S^{0}_{\text{Bender}}$ & Fichas de Bender \emph{al empezar la mano}: su \code{saldo} al repartir más lo que ya puso en ciegas (\code{saldo} $+$ \code{apuesta\_ronda} en \code{mano\_inicio}). \\
$S^{0}_i$ & Lo mismo para cada rival al que se le repartieron cartas en esa mano (aunque luego se retire o vaya all-in). \\
\end{params}
\uso{Se calcula UNA vez por mano, con las fichas del inicio, y el modo queda fijo durante toda la mano. Nunca se usan las fichas que quedan detrás en medio de la mano: un rival all-in tiene \code{saldo} $=0$ y eso no lo hace un rival de 0\,bb (la definición anterior, ambigua, mandaba a push/fold ante cualquier all-in de un rival y llegó a retirar AA contra un all-in de 100\,bb). Con ciegas fijas y 10\,000 fichas iniciales (100 bb), lo normal es el modo profundo; push/fold queda de respaldo para stacks realmente cortos.}
\end{formbox}

\subsection{Apertura preflop (nadie subió)}
\begin{formbox}{Rangos de apertura por posición}
\[
\text{abrir con } h \iff \text{rank}(h)\le \text{top}_{\text{pos}}\ \%
\qquad\qquad \text{tamaño de apertura}=2{,}5\ \mathrm{bb}
\]
\begin{center}\small\renewcommand{\arraystretch}{1.2}
\begin{tabular}{@{}lcccccc@{}}\toprule
Posición & UTG & HJ & CO & BTN & SB & BTN/SB heads-up\\ \midrule
top \% & 14{,}3 & 18{,}9 & 25{,}5 & 43{,}3 & 33{,}0 & 70\\ \bottomrule
\end{tabular}\end{center}
\begin{params}
rank$(h)$ & Posición de la clase de $h$ en el ranking de 169 clases (por equity contra una mano al azar), medida en \% acumulado de combos. \\
top$_{\text{pos}}$ & Porcentaje de manos que abre desde esa posición. Con 5 jugadores no hay UTG. \\
\end{params}
\uso{Bender nunca entra pagando la ciega (sin limp): abre subiendo o se retira. Si ya hubo una subida, decide con el EV (§1.6–§1.9) usando el rango estimado del que subió.}
\end{formbox}

\subsection{Pot odds y EV de igualar}
\begin{formbox}{Igualar (call)}
\[
e_{\min}=\frac{C}{P+C}
\qquad\qquad
EV(\text{call})=e\,(P+C)-C
\]
\begin{params}
$P$ & Pozo \textbf{antes} de igualar, incluida la apuesta del rival. \\
$C$ & Lo que cuesta igualar. \\
$e$ & Equity contra el rango del rival (§1.3). \\
\end{params}
\uso{Igualar conviene si $e>e_{\min}$, lo mismo que $EV(\text{call})>0$. Supone que no hay más apuestas después; si quedan calles, es una aproximación.}
\ej{$P=1300$, $C=900$: $e_{\min}=40{,}9\,\%$. Con $e=25{,}1\,\%$: $EV=0{,}251\cdot2200-900\approx-348$ $\Rightarrow$ retirarse.}
\end{formbox}

\subsection{Fold equity}
\begin{formbox}{Probabilidad de que todos se retiren}
\[
f=\prod_{i=1}^{n}(1-c_i)
\qquad\qquad
c(x)=\frac{0{,}65}{1+0{,}55\,x}\ \ \big(+0{,}12\ \text{si el rival ya puso fichas en esta ronda; tope } 0{,}9\big)
\]
\begin{params}
$n$ & Rivales que todavía pueden responder a la apuesta. \\
$c_i$ & Probabilidad de que el rival $i$ pague. Sin datos del rival se usa el modelo de población $c(x)$; con datos, el $c_i$ de la Capa 2 (§2.1 y §3.3). \\
$x$ & Tamaño de la apuesta en fracción del pozo (§1.8). \\
$0{,}65;\ 0{,}55;\ 0{,}12$ & Constantes del modelo de población. \emph{Provisional}. \\
\end{params}
\uso{En modo push/fold preflop se usa $c=0{,}385$ (rango de pago Nash a 10 bb, §1.10).}
\ej{$c(0{,}5)=0{,}51$, $c(1)=0{,}42$, $c(2)=0{,}31$. Dos rivales con $c=0{,}45$ y $0{,}60$: $f=0{,}55\cdot0{,}40=0{,}22$.}
\end{formbox}

\subsection{Tamaño de apuesta}
\begin{formbox}{Tamaño relativo al pozo (definición general)}
\[
x=\frac{\text{monto}-\text{apuesta actual}}{\text{pozo si pagas}}
\qquad
\text{y si nadie apostó en la ronda:}\quad x=\frac{B}{P}
\]
\begin{params}
monto & Total al que Bender lleva su apuesta en la ronda (lo que se envía en \code{monto}). \\
apuesta actual & La apuesta más alta de la ronda antes de actuar. \\
pozo si pagas & $P+C$: el pozo después de que Bender igualara. \\
\end{params}
\uso{Menú de 5 acciones: retirarse, pasar o igualar, $x\approx\tfrac12$, $x\approx1$ y all-in. Coincide con la cuenta de \code{ApostarPozo} del SDK del kit.}
\end{formbox}

\subsection{Faroles y apuestas}
\begin{formbox}{Farol puro (sin equity si le pagan)}
\[
EV(\text{farol})=f\,P-(1-f)\,B
\qquad\qquad
f^{*}=\frac{B}{P+B}
\]
\begin{params}
$f$ & Fold equity de esa apuesta (§1.7). \\
$f^{*}$ & Fold equity mínima para que un farol puro no pierda fichas. \\
\end{params}
\uso{Si $f<f^{*}$, un farol sin equity pierde. Apostar más grande baja $c$ pero sube $f^{*}$: con $B=P/2$, $f^*=33\,\%$; con $B=P$, $f^*=50\,\%$. Si el farol es una \textbf{subida} sobre otra apuesta, se cambia $B$ por $A$ (lo que agrega Bender): $f^*=A/(P+A)$.}
\end{formbox}

\begin{formbox}{EV de apostar o subir (forma general que usa el código)}
\[
EV(\text{apuesta})=f\,P+(1-f)\big[e'\,(P+A+A')-A\big]
\]
\[
\text{Caso simple (un rival, apuesta $B$ que paga completa): } A=A'=B\]
\[EV=f\,P+(1-f)\big[e'(P+2B)-B\big]
\]
\begin{params}
$A$ & Fichas que agrega Bender con esa acción. \\
$A'$ & Fichas que agrega el rival que paga: lo que le falta para igualar, con tope en su stack (si no alcanza, solo pone su all-in). Se toma el rival con más fichas en juego. \\
$e'(x)$ & Equity de Bender si le pagan \textbf{esa} apuesta: contra el rango que paga ese tamaño, no contra todo el rango. Se recalcula para cada tamaño del menú (ver abajo). \\
\end{params}
\[
R_{\text{paga}}(x)=\text{la fracción } c(x) \text{ más fuerte del rango del rival}
\qquad
e'(x)=e_R\big(R_{\text{paga}}(x)\big)
\]
\par\noindent{\small Ante una apuesta chica paga casi todo el rango; ante un all-in, solo su parte fuerte. Ejemplo: si ante un all-in $c=0{,}24$, $e'$ se calcula contra el 24\,\% más fuerte del rango, ordenado por la fuerza de cada combo en ese board (su equity contra una mano al azar, no contra la de Bender: el rival no conoce las cartas de Bender).}
\uso{Se calcula para cada tamaño del menú y se compara con $EV(\text{call})$, $EV(\text{check})=e\,P$ y $EV(\text{fold})=0$. Supone una sola respuesta y sin calles extra; no modela botes laterales.}
\ej{$P=600$, $B=300$, un rival con $c(0{,}5)=0{,}51\Rightarrow f=0{,}49$, $e'=0{,}35$: $EV=0{,}49\cdot600+0{,}51\,(0{,}35\cdot1200-300)\approx+355$.}
\par\noindent{\small\textbf{Calibración de los all-in} (opcional, interruptor \code{calibracion\_allin.activa}, \emph{apagada} por defecto). Medida con 40\,800 manos del historial local con las cartas de todos (6 bots de entrenamiento a 100\,bb; \code{pruebas/calibrar\_allin.py}): quien paga un all-in lo hace con sus manos \emph{más fuertes} (AUC $0{,}97$ en base y estadistico: el supuesto de $R_{\text{paga}}$ se cumple), pero con frecuencias muy distintas por bot: preflop pagan el $10{,}6\,\%$ (base), $44{,}2\,\%$ (estadistico) y $0\,\%$ (bandido, espia y conservador: 0 de 8\,855, 5\,033 y 10\,817 casos); postflop $22\,\%$ (base), $57\,\%$ (estadistico), $18\,\%$ (bandido), $63\,\%$ (espia) y $0\,\%$ (conservador). Solo en modo profundo, para un all-in, $c_i$ pasa a
$c_i^{\text{cal}}=\sum_j \pi_j\,\dfrac{k_j+s\,c_i}{n_j+s}$ ($n_j$ y $k_j$: casos medidos y pagos del tipo $j$; con un bot del kit reconocido, el suyo; $s=20$; lo que $\pi$ reparte entre tipos sin datos, nit y maníaco, se queda en $c_i$; no hay grupo \emph{station}: conservador no calibra a la station genérica). $R_{\text{paga}}$ sigue siendo la fracción $c_i^{\text{cal}}$ más fuerte.}
\end{formbox}

\subsection{Push/fold}
\begin{formbox}{All-in preflop con stack corto}
\[
EV(\text{push})=f\,P+(1-f)\big[e'\,(P+A+A')-A\big]\qquad\text{(la misma fórmula, con $A$ = todo el stack)}
\]
\begin{center}\small\renewcommand{\arraystretch}{1.2}
\begin{tabular}{@{}lcc@{}}\toprule
Nash a 10 bb, heads-up & SB empuja con & BB paga con\\ \midrule
Rango & top 56{,}6\,\% & top 38{,}5\,\%\\ \bottomrule
\end{tabular}\end{center}
\begin{params}
$f$ & $(1-0{,}385)^{n}$ con $n$ rivales por detrás (aproximación: con varios rivales, cada uno paga algo menos). \\
$e'$ & Equity si le pagan: contra el rango de pago, top 38{,}5\,\%. \\
\end{params}
\uso{Solo en modo push/fold ($S_{\text{ef}}\le15$). Las únicas opciones son all-in o retirarse. Los valores Nash se usan si nadie subió antes; si alguien ya subió, $e'$ es contra el rango estimado de ese rival.}
\ej{SB con 10 bb: $P=150$, $A=950$, $A'=900$, $f=0{,}615$, $e'=0{,}40$: $EV=92{,}3+0{,}385\,(800-950)\approx+35 \Rightarrow$ empuja.}
\end{formbox}

\subsection{Reglas de decisión}
\begin{formbox}{Cómo se elige la acción}
\[
a^{*}=\arg\max_{a\ \text{no descartada}} EV(a)
\qquad
\text{tipo de apuesta}=\begin{cases}\text{valor} & e'\ge 0{,}50\\ \text{semi-farol} & 0{,}25\le e'<0{,}50\\ \text{farol} & e'<0{,}25\end{cases}
\]
\begin{params}
Multiway & Con 3 o más rivales que \emph{pueden responder} a la apuesta (con fichas detrás, no all-in) y $e'<0{,}50$, se descartan las apuestas (sin valor no se farolea contra varios). Los rivales all-in siguen en el bote para la equity, pero no cuentan como respondedores. \\
Compromiso & SOLO en modo intermedio: si igualar o subir cuesta más de $\tfrac13$ del stack, se descarta: queda all-in o retirarse. Igualar un all-in ajeno, cuando eso pone todas las fichas de Bender, cuenta como la opción \emph{all-in} y no se descarta. En push/fold y en profundo la regla no se aplica nunca. \\
All-in profundo & En modo profundo, el all-in se descarta si $e'<0{,}50$ (parámetro \code{allin\_profundo\_solo\_valor}). \\
Subir sobre apuesta & Subir sobre una apuesta de otro solo con $e'\ge0{,}50$: sin semi-faroles ni faroles de re-subida (parámetro \code{solo\_valor\_sobre\_apuesta}). Medido en el paso 3: $+16{,}9\pm13{,}3$ bb/100 contra \code{base}. Se vuelve a evaluar con la Capa 2. \\
Mezcla & Si la segunda mejor acción cumple $EV_1-EV_2<0{,}02\cdot\max(P,\mathrm{bb})$, $EV_2\ge0$ y ninguna de las dos es retirarse, se elige la segunda con probabilidad $0{,}35$. Así Bender no es predecible. \\
\end{params}
\end{formbox}

\subsection{Perilla de riesgo del torneo}
\begin{formbox}{Margen de equity en calls grandes}
\[
e_{\min}^{\text{torneo}}=\frac{C}{P+C}+\delta
\qquad
\delta=\delta_{\max}\cdot g \qquad (\text{solo si } C>\tfrac13\,S_{\text{Bender}})
\]
\begin{params}
$\delta_{\max}$ & Margen máximo: 0{,}08. \emph{Provisional}. \\
$g$ & Entre $-0{,}5$ y $1$. Positivo si Bender va cómodo en el top 6 y quedan pocas rondas (pedir más equity, no arriesgar); 0 en las primeras rondas; negativo si va por debajo del corte al final (aceptar más varianza: quedar 7.º es igual que quedar último). Se define con datos en la Capa 3. \\
\end{params}
\uso{Solo afecta los calls que ponen en riesgo gran parte del stack: quedar en 0 elimina, y clasifican los 6 con más fichas.}
\end{formbox}

\subsection{Utilidad con colchón}
\begin{formbox}{Maximizar utilidad, no fichas}
\[
u(S)=\ln(S+K)
\qquad
a^{*}=\arg\max_{a\ \text{no descartada}} U(a)
\]
\begin{params}
$S$ & Fichas que le quedan \textbf{detrás} a Bender cuando acaba la mano (sin contar lo que hay en el pozo). \\
$K$ & El \emph{colchón}: cuántas fichas "regaladas" se suman a $S$. Con $K\to\infty$ la utilidad es lineal y se recupera el EV; con $K$ chico, quedar sin fichas pesa mucho. Valor inicial $10\,000$ (un stack inicial). \emph{Provisional}: se elige con la métrica del campeonato. \\
\end{params}
\[
\begin{aligned}
U(\text{retirarse}) &= u(S)\\
U(\text{pasar}) &= e\,u(S+P)+(1-e)\,u(S)\\
U(\text{igualar }C) &= e\,u(S+P)+(1-e)\,u(S-C)\\
U(\text{apostar }A) &= f\,u(S+P)+(1-f)\big[e'(x)\,u(S+P+A')+(1-e'(x))\,u(S-A)\big]
\end{aligned}
\]
\par\noindent{\small Es \textbf{el mismo árbol} de §1.6--§1.9, con $f$, $e$, $e'(x)$, $P$, $A$ y $A'$ idénticos; solo cambia cómo se valora cada hoja. En el call corto se iguala con $c_{\text{ef}}=\min(C,S)$ y lo no cubierto vuelve a su dueño (como en §1.6), y $A$ es lo que de verdad pone Bender (con tope de lo que el rival puede cubrir, como en §1.9).}
\begin{params}
Equivalente cierto & $EC(a)=(S+K)\big(e^{\,U(a)-u(S)}-1\big)$: las fichas que, ganadas con certeza, valdrían lo mismo que $U(a)$. Es $0$ para retirarse y se parece al EV cuando el riesgo es chico. Con la utilidad activa, la \textbf{mezcla} de §1.11 compara $EC_1-EC_2<0{,}02\,\max(P,\mathrm{bb})$ y $EC_2\ge0$ (en vez de $EV$). \\
Qué no cambia & Las fórmulas de EV de §1.6--§1.9 (se siguen calculando y registrando), los descartes de §1.11 y los filtros de §3.4. La perilla $\delta$ de §1.12 no entra en $U$: con la utilidad activa $g$ se queda en $0$. \\
\end{params}
\uso{Como $u$ es cóncava, poner el stack entero en una moneda al aire rinde menos utilidad que su EV en fichas. Un all-in contra un rival que paga siempre ($f=0$, $P\approx0$, $A=A'=S$) deja $\ln(2S+K)$ si gana y $\ln K$ si pierde, y retirarse deja $\ln(S+K)$: conviene solo si
$e\ge e_{\min}^{u}=\ln\!\frac{S+K}{K}\Big/\ln\!\frac{2S+K}{K}$ (en fichas bastaría $e\ge50\,\%$).}
\begin{center}\small\renewcommand{\arraystretch}{1.2}
\begin{tabular}{@{}lcccc@{}}\toprule
$K$ & 5\,000 & 10\,000 & 20\,000 & 40\,000\\ \midrule
Equity mínima de un all-in de $S=10\,000$ a cara o cruz & 68{,}3\,\% & 63{,}1\,\% & 58{,}5\,\% & 55{,}0\,\%\\ \bottomrule
\end{tabular}\end{center}
\ej{$S=10\,000$, $K=10\,000$, all-in sin fold equity con $e'=0{,}55$: $EV=0{,}55\cdot20\,000-10\,000=+1\,000$, pero $U=0{,}55\ln30\,000+0{,}45\ln10\,000\approx9{,}815<\ln20\,000\approx9{,}903=u(S)$, o sea $EC<0$: se retira (o juega una apuesta menor). Con $e'=0{,}64$ ya conviene.}
\end{formbox}

% =====================================================================
\clearpage
\section{Capa 2 · Lectura de rivales}

\subsection{Rango actualizado por cada acción}
\begin{formbox}{Reajuste del rango (Bayes con pesos)}
\[
w'(h)=\max\!\big(w(h)\,P(a\mid h),\ 0{,}02\,w(h)\big)
\qquad
\tilde w(h)=\frac{w'(h)}{\sum_{h'}w'(h')}
\qquad
c=\frac{\sum_h w(h)\,P(\text{sigue}\mid h)}{\sum_h w(h)}
\]
\begin{params}
$a$ & Acción observada del rival (sacada de \code{historial\_acciones}). \\
$P(a\mid h)$ & Verosimilitud: qué tan seguido hace $a$ quien tiene $h$. Sale de la plantilla (§2.2) de su tipo. \\
$0{,}02$ & Piso: ningún combo baja de 2\,\% de su peso, para que un showdown “imposible” no rompa el modelo. \\
$\tilde w$ & Rango normalizado (suma 1). \\
$c$ & Probabilidad de que el rival \textbf{siga} (pague o suba): sale de la misma plantilla, \textbf{antes} de que Bender apueste. La plantilla está calibrada para $x\approx\tfrac12$; para otro tamaño se escala con el modelo de población: $c(x)=c_{1/2}\cdot c_{\text{pob}}(x)/c_{\text{pob}}(\tfrac12)$. \\
\end{params}
\uso{Para $c$ se usa $P(\text{sigue}\mid h)$; para reajustar el rango después de que el rival actúa, se usa la acción exacta: $P(\text{paga}\mid h)$ si pagó, $P(\text{sube}\mid h)$ si subió. Preflop el rango arranca en §2.8. En cada calle, primero se quitan los combos bloqueados por el board.}
\ej{J$\diamondsuit$7$\clubsuit$2$\spadesuit$, 259 combos (43 fuertes, 45 medios, 171 aire), rival promedio. Antes de apostar: $c=(43\cdot.95+45\cdot.65+171\cdot.15)/259=95{,}75/259=0{,}37$. Si paga: pesos $43\cdot.55,\ 45\cdot.60,\ 171\cdot.07$ y el aire baja de 66\,\% a 19\,\% del rango.}
\end{formbox}

\subsection{Plantillas por categoría de mano}
\begin{formbox}{Verosimilitudes postflop frente a una apuesta: $P(a\mid \text{categoría},t)$ \ (\emph{provisional})}
\begin{center}\small\renewcommand{\arraystretch}{1.18}
\begin{tabular}{@{}l|ccc|ccc|ccc|ccc@{}}\toprule
& \multicolumn{3}{c|}{Fuerte} & \multicolumn{3}{c|}{Media} & \multicolumn{3}{c|}{Proyecto} & \multicolumn{3}{c}{Aire}\\
Tipo & ret & paga & sube & ret & paga & sube & ret & paga & sube & ret & paga & sube\\ \midrule
TAG (\emph{calibrado}) & .14 & .51 & .34 & .29 & .66 & .05 & .57 & .41 & .02 & .76 & .24 & .00\\
Nit & .05 & .65 & .30 & .50 & .48 & .02 & .55 & .42 & .03 & .93 & .05 & .02\\
Station & .02 & .88 & .10 & .10 & .88 & .02 & .10 & .85 & .05 & .30 & .68 & .02\\
Maníaco & .02 & .38 & .60 & .15 & .55 & .30 & .10 & .45 & .45 & .40 & .25 & .35\\
Aleatorio & $\tfrac13$ & $\tfrac13$ & $\tfrac13$ & $\tfrac13$ & $\tfrac13$ & $\tfrac13$ & $\tfrac13$ & $\tfrac13$ & $\tfrac13$ & $\tfrac13$ & $\tfrac13$ & $\tfrac13$\\ \bottomrule
\end{tabular}\end{center}
\begin{params}
Fuerte & Top pair o mejor: overpair, dos pares, trío, escalera, color… \\
Media & Pareja menor: pareja de bolsillo bajo la carta alta del board, segunda o tercera pareja. \\
Proyecto & Sin pareja pero con 8 o más outs (proyecto de color o de escalera abierta). \\
Aire & Todo lo demás (un gutshot cuenta como aire). \\
ret / paga / sube & Se retira / iguala / sube. Cada trío suma 1. “Sigue” = paga + sube. \\
\end{params}
\uso{Cuando nadie apostó, se usa $P(\text{apuesta}\mid\text{cat})$: TAG (\emph{calibrado}) apuesta con fuerte .90, media .54, proyecto .22, aire .19 (pasar = 1 menos eso). Los otros tipos parten de la fila provisional del rival promedio (fuerte .65, media .35, proyecto .45, aire .25) y se multiplica por tipo: nit $\times0{,}7$, station $\times0{,}6$, maníaco $\times1{,}5$ (tope .95); aleatorio .50 en todo. TAG se midió con las cartas de todos los jugadores en 4\,800 manos de \code{base}, \code{estadistico}, \code{bandido} y \code{espia}, con \emph{shrinkage} $s=20$ por fila hacia la fila provisional. El ejemplo de §2.1 usa la fila provisional del rival promedio (“sigue” .95, .65, .15). \textbf{Station, nit, maníaco y aleatorio son plantillas genéricas: un bot concreto que se comporte distinto se modela con su propia plantilla (§3.4), no cambiando la genérica.}}
\end{formbox}

\subsection{Bayes sobre el tipo de rival}
\begin{formbox}{Posterior del tipo}
\[
P(t\mid a)=\frac{P(a\mid t)\,P(t)}{\sum_{t'}P(a\mid t')\,P(t')}
\qquad\qquad
\log\pi_t \leftarrow \lambda\log\pi_t+\log P(a\mid t)
\]
\begin{params}
$P(t)$ & Prior de población: nit .06, TAG .53, station .18, maníaco .06, aleatorio .17. \emph{Calibrado} con los bots del kit; se revisa con el servidor de pruebas. \\
$P(a\mid t)$ & Verosimilitud de la acción para el tipo (tabla de abajo preflop; postflop, la plantilla de §2.2). \\
$\pi_t$ & Posterior guardado por ID del rival, en disco. \\
$\lambda$ & Olvido (§2.6). Con $\lambda=1$ es Bayes puro. \\
\end{params}
\begin{center}\small\renewcommand{\arraystretch}{1.18}
\begin{tabular}{@{}lccccc@{}}\toprule
Preflop (TAG, station y aleatorio \emph{calibrados}) & Nit & TAG & Station & Maníaco & Aleatorio\\ \midrule
Se retira (F) & .88 & .65 & .06 & .35 & .30\\
Entra pagando (C) & .04 & .26 & .92 & .15 & .14\\
Entra subiendo (R) & .08 & .09 & .02 & .50 & .56\\ \bottomrule
\end{tabular}\end{center}
\uso{Se trabaja en logaritmos (evita \emph{underflow}). Para normalizar: restar el máximo, exponenciar y dividir por la suma. Piso: ningún $\pi_t$ baja de 0{,}01.}
\ej{Prior + “entra pagando”: station pasa de 20\,\% a 50{,}6\,\% en una mano.}
\end{formbox}

\subsection{Mezcla de plantillas}
\begin{formbox}{Verosimilitud que usa Bender para el rango}
\[
P(a\mid h)=\sum_{t}\pi_t\;P(a\mid h,t)
\]
\begin{params}
$\pi_t$ & Creencia actual sobre el tipo del rival (§2.3). \\
$P(a\mid h,t)$ & Plantilla del tipo $t$ para la categoría de $h$ (§2.2). \\
\end{params}
\uso{No se apuesta por un solo tipo: el rango y el $c$ salen de la mezcla. Si el rival es un bot del kit reconocido (§2.7), se usa su lógica en lugar de la mezcla.}
\end{formbox}

\subsection{Estadísticas con prior Beta (shrinkage)}
\begin{formbox}{Estimar una frecuencia con pocas manos}
\[
\alpha_0=s\,m,\quad \beta_0=s\,(1-m)
\qquad
\text{tras $k$ síes en $n$:}\quad \alpha=\alpha_0+k,\ \ \beta=\beta_0+n-k
\]
\[
\hat p=\frac{\alpha}{\alpha+\beta}=\frac{k+s\,m}{n+s}
\qquad\qquad
\sigma=\sqrt{\frac{\hat p\,(1-\hat p)}{\alpha+\beta+1}}
\]
\begin{params}
$k,\ n$ & Veces que pasó lo que se mide / oportunidades observadas. \\
$m$ & Promedio de la población (se mide en el servidor de pruebas). \emph{Calibrado}: VPIP 0{,}455, PFR 0{,}143, sigue ante apuesta chica 0{,}692, sigue ante grande 0{,}545. \\
$s$ & Fuerza del prior, en manos imaginarias: 20 para VPIP y PFR. \emph{Provisional}. \\
$\hat p$ & Estimado de la frecuencia. \\
$\sigma$ & Incertidumbre del estimado: decide cuánto explotar (§3.2). \\
\end{params}
\uso{Estadísticas por ID: VPIP, PFR, “sigue (paga o sube) ante una apuesta chica” ($x\le0{,}6$) y “sigue ante una apuesta grande” ($x>0{,}6$). Cada una es un par $(\alpha,\beta)$.}
\ej{$k=7$, $n=10$, $m=0{,}3$, $s=20$: $\hat p=13/30=43{,}3\,\%$ (crudo: 70\,\%), $\sigma=8{,}9\,\%$.}
\end{formbox}

\subsection{Olvido}
\begin{formbox}{Adaptarse a rivales que cambian}
\[
\alpha\leftarrow\lambda\,\alpha,\quad \beta\leftarrow\lambda\,\beta \ \ \text{(primero)} \qquad\text{luego se suma la observación nueva}
\qquad
\text{memoria}=\frac{1}{1-\lambda}
\]
\begin{params}
$\lambda$ & Factor de olvido: 0{,}98 (memoria de unas 50 manos). \emph{Provisional}, rango razonable 0{,}97–0{,}99. \\
\end{params}
\uso{Se aplica igual a $\log\pi_t$ (§2.3). Es la misma idea que el horizonte $1/(1-\gamma)$ de los MDP. Con olvido, $\sigma$ nunca llega a 0, lo que frena la explotación.}
\end{formbox}

\subsection{Huella de los montos (bots del kit)}
\begin{formbox}{Reconocer un bot copiado del kit}
\[
z=\frac{\text{monto}-\text{apuesta actual}}{\text{pozo si pagas}}
\qquad\qquad
P(\text{bot}\mid z_1,\dots,z_n)\propto P(\text{bot})\prod_{i}P(z_i\mid \text{bot})
\]
\begin{params}
$z$ & Tamaño de cada apuesta postflop de un rival (misma cuenta que $x$). Cae en una caja si está a $\pm0{,}03$ de .33, .4, .5, .6, .66, .75, .9 o 1.0; si no, en “otro”. No se cuentan los all-in ni las apuestas recortadas por el stack. \\
$P(\text{bot})$ & 0{,}5 para “propio” (bot escrito desde cero) y 0{,}1 para cada bot del kit. \\
$P(z\mid\text{bot})$ & Tabla de abajo, \emph{calibrada} midiendo cada bot en 90 mesas. $^{\dagger}$\code{novato} es de reserva y no se midió: conserva la estimación leída del código. \code{conservador} nunca apuesta y no tiene huella. “Propio” reparte $1/9$ en cada caja. \\
\end{params}
\begin{center}\small\renewcommand{\arraystretch}{1.15}
\begin{tabular}{@{}lccccccccc@{}}\toprule
Bot & .33 & .4 & .5 & .6 & .66 & .75 & .9 & 1.0 & otro\\ \midrule
base & .01 & .77 & .06 & .01 & .01 & .11 & .01 & .01 & .01\\
novato\,$^{\dagger}$ & .20 & .01 & .44 & .01 & .01 & .29 & .01 & .01 & .02\\
espia & .01 & .01 & .33 & .47 & .01 & .01 & .10 & .06 & .01\\
bandido & .01 & .01 & .34 & .01 & .57 & .01 & .01 & .04 & .01\\
estadistico & .01 & .01 & .18 & .11 & .12 & .15 & .07 & .03 & .32\\
aleatorio & .01 & .03 & .02 & .03 & .02 & .01 & .01 & .02 & .84\\ \bottomrule
\end{tabular}\end{center}
\uso{Si $P(\text{bot})\ge0{,}9$, Bender usa la contraestrategia de ese bot (§3.4). Si el equipo cambió los tamaños, el posterior se va a “propio” y no se rompe nada.}
\ej{0{,}75, 0{,}4, 0{,}75 $\Rightarrow$ base con 93\,\%.}
\end{formbox}

\subsection{Rango preflop inicial desde las estadísticas}
\begin{formbox}{Ancho del rango cuando el rival entra}
\[
\text{rango si sube}=\text{top}\big(100\,\hat p_{\text{PFR}}\big)\,\%
\qquad
w_{\text{paga}}(h)=\begin{cases}1 & h\in\text{top}(100\,\hat p_{\text{VPIP}})\setminus\text{top}(100\,\hat p_{\text{PFR}})\\ 0{,}3 & h\in\text{top}(100\,\hat p_{\text{PFR}})\\ 0 & \text{resto}\end{cases}
\]
\begin{params}
$\hat p_{\text{VPIP}}$ & Estimado de qué \% de manos juega (§2.5). \\
$\hat p_{\text{PFR}}$ & Estimado de qué \% de manos sube preflop. \\
\end{params}
\uso{Quien solo paga suele no tener sus mejores manos (con esas habría subido), así que el top del PFR entra con peso 0{,}3 (\emph{provisional}: a veces paga con una mano fuerte para tender una trampa). Es el punto de partida del rango antes del flop. Quien sube preflop parte del top \% de su posición (§1.5), y cada re-subida multiplica el ancho por $0{,}45$ (\emph{provisional}). Mientras $\rho$ sea bajo (§3.2), se mezcla con el top \% de población de esa posición.}
\par\noindent{\small\textbf{Rango de quien empuja all-in} (misma calibración opcional, apagada por defecto). Medido en las mismas 40\,800 manos, preflop, el 90\,\% de los all-in de cada bot sale del top: base $9{,}8\,\%$, espia $3{,}2\,\%$, estadistico $2{,}7\,\%$, bandido $1{,}8\,\%$ (aleatorio, $90\,\%$: empuja con cualquier cosa). Con la calibración activa y en modo profundo, el rango de un rival que hizo all-in preflop se recorta al top $T\,\%$ con $T=\sum_j\pi_j\,\dfrac{n_j\,q^{90}_j+s\cdot100}{n_j+s}$ (shrinkage hacia $100$, o sea sin recorte). Postflop, la probabilidad de \emph{empujar} dada la categoría de la mano (§2.2) también se mide: con mano fuerte y apuesta enfrentada empujan el $8{,}7\,\%$ (base), $34{,}8\,\%$ (estadistico), $21{,}9\,\%$ (espia) y $6{,}8\,\%$ (bandido); con mano media o menos, casi nunca. Cada celda se encoge hacia la plantilla actual con $s=20$ y entra en el reajuste de §2.1 como $P(\text{all-in}\mid\text{categoría})$.}
\end{formbox}

% =====================================================================
\section{Explotación · de la lectura a la decisión}

\begin{notebox}\small
\textbf{Principio.} Explotar no es agregar reglas nuevas: es meter los números \textbf{del rival} ($c_i$ y su rango) en \textbf{las mismas fórmulas} de la Capa 1. Si el rival paga todo, su $c$ es alto: $f$ baja, los faroles pierden EV y las apuestas de valor grandes ganan EV, sin que nadie lo programe a mano. Lo único nuevo es \textbf{cuánto} confiar en esos números.
\end{notebox}

\subsection{Peso de explotación según la confianza}
\begin{formbox}{Cuánto se desvía Bender de la población}
\[
\rho_i=\rho_{\max}\cdot\max\!\left(0,\ 1-\frac{\kappa\,\sigma_i}{|\hat p_i-m|}\right)
\qquad (\rho_i=0 \text{ si } n_i<n_{\min} \text{ o } \hat p_i=m)
\]
\begin{params}
$\hat p_i,\ \sigma_i$ & Estimado e incertidumbre de la estadística del rival $i$ (§2.5). \\
$m$ & Valor de población de esa estadística. \\
$|\hat p_i-m|$ & Cuánto se aleja el rival de lo normal. \\
$\kappa$ & Exigencia: cuántos $\sigma$ tiene que separarse para explotar. 1{,}0 \emph{provisional}. \\
$\rho_{\max}$ & Tope de explotación: 0{,}8 (siempre queda algo de juego sólido). \\
$n_{\min}$ & Mínimo de observaciones: 8. \\
\end{params}
\uso{Si la diferencia con la población es menor que $\kappa\sigma$, no se explota ($\rho=0$). A medida que $\sigma$ se achica, $\rho$ sube hasta $\rho_{\max}$. Es conservador a propósito: $\hat p$ ya viene encogido hacia $m$ (§2.5) y $\rho$ vuelve a frenar.}
\ej{VPIP con $m=0{,}30$. Con 10 manos ($\hat p=0{,}433$, $\sigma=0{,}089$): $\rho=0{,}8\,(1-0{,}089/0{,}133)=0{,}26$. Con 85 manos ($\hat p=0{,}552$, $\sigma=0{,}048$): $\rho=0{,}8\,(1-0{,}048/0{,}252)=0{,}65$.}
\end{formbox}

\subsection{Mezcla de parámetros}
\begin{formbox}{Parámetro que entra a las fórmulas}
\[
\theta_i=(1-\rho_i)\,\theta_{\text{base}}+\rho_i\,\hat\theta_i
\]
\begin{params}
$\theta_{\text{base}}$ & Valor del modelo. Para $c_i$: el $c$ que sale del rango y las plantillas (§2.1), que ya usa el tipo y el board; sin rango, $c(x)$ de §1.7. Para el rango preflop: el top \% de población de esa posición. \\
$\hat\theta_i$ & Valor medido del rival: $\hat p$ de “sigue ante apuesta chica/grande” (§2.5), o el ancho de §2.8. \\
$\rho_i$ & Peso de explotación de esa estadística (§3.2). \\
\end{params}
\uso{Se aplica a dos cosas: el $c_i$ de cada rival (según el tamaño $x$, caja chica o grande) y el ancho de su rango preflop. Lo demás sale solo de las fórmulas de la Capa 1.}
\end{formbox}

\subsection{Ajustes por tipo (solo con certeza alta)}
\begin{formbox}{Reglas extra cuando $\pi_t\ge\tau$}
\begin{params}
Nit & Robarle las ciegas seguido; respetar sus subidas grandes (su rango ya es estrecho por §2.1). \\
TAG & Juego sólido de Capa 1, sin ajustes extra. \\
Station & Apuestas de valor con $x\approx1$ incluso con manos medias-fuertes. El filtro “no farolearle nunca” está apagado por defecto (\code{filtros\_tipo}): el EV con el $c_i$ medido ya decide si un farol rinde. \\
Maníaco & No farolearle; pagar más liviano (su rango de subida es ancho) y dejar que apueste él. \\
Aleatorio & Apuestas de valor finas y frecuentes; no darle crédito a sus subidas. \\
Bot del kit & Con $P(\text{bot})\ge0{,}9$, su plantilla medida reemplaza a la mezcla de tipos. Para \code{base}, además, el rango preflop sale de su código: sube con Chen $\ge9$ y paga con $5\le$ Chen $<9$. \code{conservador} se reconoce por comportamiento (nunca apuesta). \\
$\tau$ & Umbral de certeza: 0{,}7. \emph{Provisional}. \\
\end{params}
\uso{Estas reglas solo filtran el menú (por ejemplo, descartar faroles contra station). Los EV siguen saliendo de las fórmulas.}
\end{formbox}

% =====================================================================
\section{Parámetros ajustables (los afina la Capa 3)}

\begin{center}\small\renewcommand{\arraystretch}{1.2}
\begin{tabularx}{\linewidth}{@{}>{\raggedright\arraybackslash}p{3.6cm}>{\centering\arraybackslash}p{3.2cm}>{\raggedright\arraybackslash}X@{}}
\toprule
{\labelfont Parámetro} & {\labelfont Valor inicial} & {\labelfont Dónde se usa}\\ \midrule
Reloj adaptativo & $N_{\min}$ 2000 / $\varepsilon_{\text{obj}}$ 0{,}005 / $k$ 3 / tope 1{,}5 s / $\eta$ 0{,}5 & Monte Carlo (§1.2) \\
Umbrales de modo & 15 / 25 bb & Modo de juego (§1.4) \\
Tamaño de apertura & 2{,}5 bb & Preflop (§1.5) \\
Constantes de $c(x)$ & 0{,}65 / 0{,}55 / +0{,}12 & Fold equity de población (§1.7) \\
Umbrales de tipo de apuesta & 0{,}50 / 0{,}25 & Valor, semi-farol, farol (§1.11) \\
Regla de compromiso & $\tfrac13$ del stack & Modo intermedio (§1.11) \\
Mezcla & 2\,\% del pozo / 0{,}35 & Aleatorizar empates de EV (§1.11) \\
Reglas de solo valor & ambas activas & All-in profundo y subir sobre apuesta (§1.11) \\
Estrechar por re-subida & $\times0{,}45$ & Rango preflop (§2.8) \\
$\delta_{\max}$ & 0{,}08 & Perilla de riesgo (§1.12) \\
$K$ (colchón) & 10\,000 (utilidad \emph{encendida} por defecto desde \code{v1.1-riesgo}) & Utilidad $\ln(S+K)$ (§1.13) \\
Calibración de all-in & $s=20$ (apagada por defecto) & $R_{\text{paga}}$ y rango de quien empuja (§1.9, §2.8) \\
Piso de rango & 0{,}02 & Reajuste del rango (§2.1) \\
Plantillas & tabla §2.2 & Rango y $c$ (§2.1) \\
Prior de tipos & .06/.53/.18/.06/.17 & Bayes (§2.3) \\
$m$, $s$ & medido / 20 & Shrinkage (§2.5) \\
$\lambda$ & 0{,}98 & Olvido (§2.6) \\
Umbral de huella & 0{,}9 & Bot reconocido (§2.7) \\
$\kappa$, $\rho_{\max}$, $n_{\min}$ & 1{,}0 / 0{,}8 / 8 & Peso de explotación (§3.2) \\
$\tau$ & 0{,}7 & Ajustes por tipo (§3.4) \\
\bottomrule
\end{tabularx}
\end{center}

\end{document}
